% 原书分章PDF第11—14页：第30.5节全部11段，I/II/III三联结构完整。
\section{真知识的三联画}

% S30.5.P01
在思想史中，对真知识或可靠知识的追求，与对所用方法及其（暂时或最终）发现的验证，始终紧密相连。对于上述努力最终形成的各种书面表达，只要作善意的理解，就可以辨认出一种分层结构，并用一个简单的“架构”来描述。Starmans（2017d）认为，追求真理同样具有分层结构：经验科学（I）、研究方法论（II），以及认识论／科学哲学（III）。

% S30.5.P02
最具体、最特定的层次，位于各个专门学科（I）之内；为方便起见，这里把经验科学作为范例。某一科学领域的实践者拥有\emph{领域知识}（domain knowledge），并试图解决该领域的\emph{知识问题}（knowledge problems）。严格地说，他们在“生产”知识。为此，研究者描述实在的某些方面，并为之建立模型、寻找解释。理想情况下，这些解释构成某种尽可能普遍、理想上具有统一性的理论的一部分，而这种理论在说明和预测方面都能取得成功。为实现这一点，研究者作观察、收集数据、进行实验、研究因果关系、尝试表述（自然）规律、建立模型，并检验假说。

% S30.5.P03
通常，某一特定知识领域中的大多数知识问题，无论属于法律、医学、商学、自然科学还是社会科学，都具有因果性。它们预设经验领域背后的实在中存在原因—结果关系，或者至少确立了在上述理论中表达出来、或潜藏于理论之下的因果内涵。研究问题具有因果性，研究问句具有因果性，论证模式具有因果性；实际上，无论出自外行还是专家，用来表达这些知识问题的整套自然语言，也具有因果性。无论面对的是日常语言话语、实践问题，还是高度理论化的问题，都很难想象一个不具有因果性的研究问题。因此，因果性话语是无处不在、内在其中的。如果人们想要认识、解释、建模或理解某种事物，因果性就会浮现，并且似乎与领域知识、实践水平和现有理论、问题的性质，以及该学科的目标交织在一起。

% S30.5.P04
在层次I，因果性不仅显而易见、内在其中，而且在某种意义上还具有语境敏感性（context-sensitive），取决于相关领域的习惯、特征和独特性，无论该领域是医学、法律、自然科学、社会科学，还是商学。实际上，正如随后两节将要勾勒的那样，这种语境敏感性在哲学史中因果性概念的发展过程中也一直可见。无论如何，上述这幅或多或少存在问题的科学“标准图景”，有一个根本方面：经验与理论之间微妙的关系，或被推测的联系与互动。一边是观察、数据、“特殊的”与具体的东西，另一边则是规律、假说、“一般的”与抽象的东西。

% S30.5.P05
认识论在很大程度上正是凭借二者之间被假定或被体验到的张力场而存在。一般的知识获取，尤其是与之相关的推理或推断，始终需要推理模式的倒转或逆向。其“方向”在特殊与一般之间、具体与抽象之间、有限与无限之间，以及数据与假说／理论之间往复变换；而最重要的是，在原因与结果之间往复变换。事实上，Gelman 和 Shalizi（2013）试图从认识论角度加以证成的贝叶斯方法，往往因能够方便而有说服力地对这一侧面建模而受到赞赏。毋庸置疑，关于因果性的这些直觉和洞见，不仅体现在经验科学中，而且在工程学科或设计科学（design science）中更为重要。在这些领域，设计解决方案是为了应对行动问题，而不是知识问题；科学过程由设计循环引导，而不是由经验循环引导（Wieringa，2014）。

% S30.5.P06
开展这类研究时，研究者可以依赖第二个层次，即方法论层次（II）。它更加抽象、较少依赖具体领域，提供形式化的方法与表征、技术与工具，帮助研究者实际开展研究——包括上述推理模式的逆向——并验证结果。可靠而有效的结论，要求弥合经验与理论之间被认为存在的鸿沟。有时，这些方法具有一般性，例如统计方法、数据分析算法（机器学习、数据挖掘），或研究设计方法。有时，则需要特定（学科）的方法和技术，例如测量工具、实验环境、仿真模型或验证工具。

% S30.5.P07
对于这一层次的因果性问题，目前只能说，方法论研究者有责任帮助研究者回答其研究问句，包括其中的因果内涵、暗示或话语；（自然）科学的进步，过去与现在都取决于这一联盟在多大程度上成功回答、重新表述或形式化了因果问题，或者可能回避、绕开了它们。例如，在（实验）方法论中，一系列统计技术正发挥着这样的作用，包括（典型）相关、偏相关、列联表分析与关联度量、方差分析、Logistic与线性回归分析、路径模型，以及结构比较模型（structural comparison models）。此外，因子分析、主成分分析（principle component analysis）、判别分析，甚至时间序列预测与生存分析，也部分带有因果内涵，或至少允许作因果解释。直到今天，这些方法仍主导着实验研究和观察性研究，并构成许多经验学科中因果陈述的基础。

% S30.5.P08
严格从实用与科学主义（scientism）的意义上说，对知识基础及其证成的追寻，可以被认为已经完成。然而，科学史表明，人们通常还会诉诸第三个层次。认识论／一般科学哲学（III）补全了这一架构。尤其在那些刚刚经历过一次或多次基础危机、并已达到某种成熟程度的学科中，这一点十分明显。研究者明确地寻求与认识论的重大主题、已识别的科学观念，以及受到珍视或被拆解的知识理想建立联系，并以之为锚。这类学科的哲学因此会引入如下主题：有效而可靠的知识、知识获取、真理、实在、陈述、因果性、解释与自然规律、科学理论的结构、理性与科学进步，以及各门科学的统一化或统一性。

% S30.5.P09
如今，科学哲学中综合性、普遍性理论的时代已经过去；除一般科学哲学之外，特殊科学哲学也获得了越来越多空间。后者关注特定问题，例如数学、物理学、生物学、经济学、社会科学、神经科学或计算机科学中的问题。人们显然希望实现的互动，至少可以从两个核心问题出发：在多大程度上，我们能够依据一般科学哲学所提供的普遍洞见、理论和概念，理解某一特殊科学领域内部的发展？反过来，某一特殊科学领域，又能够为一般知识研究与科学研究中现有理论的进一步发展作出什么贡献？

% S30.5.P10
本文只能限于指出，无论如何，这个简单架构至少包含三种关系：I与II之间、II与III之间，以及I与III之间的关系。思想史中的许多争论，以及当代科学与哲学关系中的瓶颈，都可以从这些关系推导出来；许多问题也可以从这一视角得到理解和说明（Starmans，2017b,d）。例如，本节表明，物理学中I与III之间的关系问题重重；Gelman 和 Shalizi（2013）的哲学转向，主要涉及II与III之间的关系；而正如下面几节将要看到的那样，当前关于如何表征因果性的争论，分别涉及I与II之间或I与III之间的关系。这里有许多方面只能略而不论，但此时应当强调，Gelman 和 Shalizi（2013）试图联系II—III的努力值得赞赏；不过，爱是相互的，必须来自双方。

% S30.5.P11
认识论／一般科学哲学的一大问题是：就统计学而言，认识论固守发现语境与证成语境之间的古老区分，并不试图把握或讨论统计学的发展，贝叶斯统计中的少数例外除外。这一点似乎也内在于因果性文献之中：这些文献沿上述界线被大幅分割，令人遗憾，而且几乎没有相互启发。近期的一项研究（Phyllis 和 Russo，2016）对研究实践表现出一定兴趣，无疑是弥合鸿沟的一次有意思的尝试；但是，统计科学——包括 TMLE——与算法数据分析中关于因果性的近期发展，仍在很大程度上受到忽视。下一节将勾勒这一问题的若干方面。
